%% ma: L10-B14-0065
%% lop: 10
%% chuong: 5
%% bai: 14
%% dang: 10.14.10
%% loai: TN4
%% muc: TH
%% dap_an: "C"
%% nguon: "chuyen-de-cac-so-dac-trung-cua-mau-so-lieu-khong-ghep-nhom-toan-10-chuong-trinh-moi.pdf (GV. Phan Nhật Linh) (P05), Bài 03 Dạng 2 Phần I Câu 9"
%% kien_thuc: "Lớp 10 Bài 13 (số trung bình, trung vị, tứ phân vị, mốt), Bài 14 (khoảng biến thiên, khoảng tứ phân vị, phương sai, độ lệch chuẩn, giá trị bất thường)"
%% trang_thai: chinh-thuc
%% ghi_chu: "Tự duyệt theo yêu cầu giáo viên 22/10/2026."
\begin{ex}
Chọn khẳng định đúng trong các khẳng định sau.
\choice{Các số đo độ phân tán có thể nhỏ hơn $0$}{Khoảng tứ phân vị có sử dụng thông tin của giá trị lớn nhất và nhỏ nhất}{\True Khoảng biến thiên chỉ sử dụng thông tin của giá trị lớn nhất và nhỏ nhất, bỏ qua thông tin của các giá trị còn lại}{Nếu các giá trị của mẫu càng tập trung quanh giá trị trung bình thì độ lệch chuẩn càng lớn}
\loigiai{A sai (các số đo độ phân tán không âm); B sai ($\Delta_Q$ chỉ dùng $50\%$ số liệu ở giữa); D sai (càng tập trung thì độ lệch chuẩn càng nhỏ). C đúng theo định nghĩa $R=x_{\max}-x_{\min}$. Chọn C.}
\end{ex}
